% -----------------------------------------------------------
\section{Faith}
% -----------------------------------------------------------
	\label{FAITH}
	
% % ----------------------
% \subsection{See also}
% % ----------------------

% % ----------------------
% \subsection{1828 American Dictionary of the English Language} % *1828
% % ----------------------
	% \label{FAITH-1828}

	% \begin{compactitem}
		% \item
	% \end{compactitem}

% % ----------------------
% \subsection{Oxford English Dictionary} % *OED
% % ----------------------
	% \label{FAITH-OED}

	% \begin{compactitem}
		% \item[]
	% \end{compactitem}

% ----------------------
\subsection{Scriptural Usage} % *SCRIPTURES
% ----------------------
	\label{FAITH-SCRIPTURES}

	% % -----
	% \subsubsection{Old Testament} % *OT
	% % -----
		% \label{FAITH-OT}
	
		% % --
		% \paragraph{KJV}
		% % --
			% \label{FAITH-OT-KJV}
	
			% \begin{tabular}{ll}
				% Total occurrences in the OT & 
			% \end{tabular}
			
			% \begin{compactdesc}
				% \item[]
			% \end{compactdesc}
			
		% % --
		% \paragraph{Hebrew Bible}
		% % --
			% \label{FAITH-HEBREW}
		
			% %
			% \subparagraph{\texorpdfstring{\he{sample word}}{}}
			% %
				% \label{PAGA} % whatever is in step bible without periods
			
				% \begin{compactdesc}
					% \item[HALOT , PDF ] \textbf{}
						% \begin{compactitem}
							% \item[1]
						% \end{compactitem}
					% \item[BDB , PDF ] \textbf{}
						% \begin{compactitem}
							% \item[1]
						% \end{compactitem}
					% \item[DCH PDF ] \textbf{}
						% \begin{compactitem}
							% \item[1]
						% \end{compactitem}
				% \end{compactdesc}
				
				% \begin{compactdesc}
					% \item[Some OT uses] \textbf{}
						% \begin{compactdesc}
							% \item[]
						% \end{compactdesc}
				% \end{compactdesc}
			
		% % --
		% \paragraph{Septuagint}
		% % --
			% \label{FAITH-LXX}
		
			% %
			% \subparagraph{\texorpdfstring{\gr{sample word}}{}}
			% %
		
	% -----
	\subsubsection{New Testament} % *NT
	% -----
		\label{FAITH-NT}
	
		% --
		\paragraph{KJV}
		% --
			\label{FAITH-NT-KJV}
			
	
			% \begin{tabular}{ll}
				% Total occurrences in the NT & 
			% \end{tabular}
			
			\begin{compactdesc}
				\item[{\hyperref[JOHN-20-24]{John 20:24--29}}] \phantom{.}
					\begin{compactitem}
						\item Compare \hyperref[3NEPHI-11-15]{3 Nephi 11:15}.
					\end{compactitem}
				\item[{\hyperref[ROMANS-8-24]{Romans 8:24--25}}]
			\end{compactdesc}
			
		% % --
		% \paragraph{Greek New Testament} % *GREEK
		% % --
			% \label{FAITH-GREEK}
		
			% %
			% \subparagraph{\texorpdfstring{\gr{sample word}}{}}
			% %
				% \label{KATALLAGE}
			
				% \begin{compactdesc}
					% \item[BDAG , PDF ] \textbf{}
						% \begin{compactitem}
							% \item 
						% \end{compactitem}
					% \item[LSJ , PDF ] \textbf{}
						% \begin{compactitem}
							% \item 
						% \end{compactitem}
				% \end{compactdesc}
				
				% \begin{compactdesc}
					% \item[Some NT uses] \textbf{}
						% \begin{compactdesc}
							% \item[]
						% \end{compactdesc}
				% \end{compactdesc}
		
	% -----
	\subsubsection{Book of Mormon} % *BOM
	% -----
		\label{FAITH-BOM}
	
		% \begin{tabular}{ll}
			% Total occurrences in the BOM & 
		% \end{tabular}
		
		\begin{compactdesc}
			\item[{\hyperref[1NEPHI-7-12]{1 Nephi 7:12}}]
			\item[{\hyperref[2NEPHI-26-13]{2 Nephi 26:13}}]
			\item[{\hyperref[2NEPHI-27-23]{2 Nephi 27:23}}]
			\item[{\hyperref[ALMA-32]{Alma 32}}]
			\item[{\hyperref[MORMON-3-12]{Mormon 3:12}}] \phantom{.}
				\begin{compactitem}
					\item Compare \hyperref[MORMON-5-2]{Mormon 5:2}.
				\end{compactitem}
			\item[{\hyperref[MORMON-8-24]{Mormon 8:24}}]
			\item[{\hyperref[MORMON-9-20]{Mormon 9:20--21}}]
			\item[{\hyperref[ETHER-12]{Ether 12}}]
			\item[{\hyperref[MORONI-7]{Moroni 7}}] \phantom{.}
				\begin{compactitem}
					\item Various verses in this chapter talk about it, and they're important.
				\end{compactitem} 
			\item[{\hyperref[MORONI-10-7]{Moroni 10:7}}]
			\item[{\hyperref[MORONI-10-19]{Moroni 10:19--26}}]
		\end{compactdesc}
		
	% -----
	\subsubsection{Doctrine and Covenants} % *DC
	% -----
		\label{FAITH-DC}
	
		% \begin{tabular}{ll}
			% Total occurrences in the DC & 
		% \end{tabular}
		
		\begin{compactdesc}
			\item[{\hyperref[DC10-48]{DC 10:48--49}}] \phantom{.}
				\begin{compactitem}
					\item Explaining a particular instance of faith by specifying some content with a desire.
				\end{compactitem}
			\item[{\hyperref[DC11-17]{DC 11:17}}] \phantom{.}
				\begin{compactitem}
					\item You can read this verse as saying that faith has a desire-based component as well. So not just a knowledge one. Similar to DC 10:48--49 above.
				\end{compactitem}
			\item[{\hyperref[DC42-43]{DC 42:43--52}}] \phantom{.}
				\begin{compactitem}
					\item Notice that in both verse 43 and verse 52, the Lord draws a distinction between having faith (or maybe having a sufficient amount of faith) and believing in him.
				\end{compactitem}
			\item[{\hyperref[DC46-13]{DC 46:13--14}}]
			\item[{\hyperref[DC46-19]{DC 46:19--20}}]
		\end{compactdesc}
		
	% % -----
	% \subsubsection{Pearl of Great Price} % *PGP
	% % -----
		% \label{FAITH-PGP}
	
		% \begin{tabular}{ll}
			% Total occurrences in the PGP & 
		% \end{tabular}
		
		% \begin{compactdesc}
			% \item[]
		% \end{compactdesc}
		
% % ----------------------
% \subsection{Key Phrases} % *KEYPHRASES *PHRASES
% % ----------------------
	% \label{FAITH-PHRASES}

	% \begin{compactdesc}
		% \item[]
	% \end{compactdesc}
	
% % ----------------------
% \subsection{Bible Dictionaries} % *DICTIONARIES
% % ----------------------
	% \label{FAITH-BD}

% % ----------------------
% \subsection{Restoration Usage} % *RESTORATION
% % ----------------------
	% \label{FAITH-RESTORATION}

	% % -----
	% \subsubsection{Joseph Smith} % *JS
	% % -----
		% \label{FAITH-JS}
	
		% \begin{compactdesc}
			% \item[DATE/SOURCE]
		% \end{compactdesc}
	
	% % -----
	% \subsubsection{Brigham Young} % *BY
	% % -----
		% \label{FAITH-BY}
	
		% \begin{compactdesc}
			% \item[DATE/SOURCE]
		% \end{compactdesc}
	
% ----------------------
\subsection{Commentary and Studies} % *COMMENTARY *STUDIES
% ----------------------
	\label{FAITH-COMMENTARY}

	% -----
	\subsubsection{Key Points}
	% -----
		\label{FAITH-KEYS}
	
		\begin{compactitem}
			\item Possibly see faith in a more Aristotelian way, as a habit or pattern of behavior? See my note at \hyperref[2NEPHI-9-23]{2NEPHI-9-23}, and also see the thought below from Roger Johnson.
		\end{compactitem}
		
	% -----
	\subsubsection{Faith, hope, and charity}
	% -----
		\label{FAITH-hope-charity}
		
		A list of relevant scriptures:
		
			\begin{compactitem}
				\item \hyperref[ROMANS-4-18]{Romans 4:18}
				\item \hyperref[DC18-19]{DC 18:19}
				\item \hyperref[DC67-3]{DC 67:3}
			\end{compactitem}
			
	% -----
	\subsubsection{Faith and works}
	% -----
		\label{FAITH-works}
		
		See the following places:
		
			\begin{compactitem}
				\item Notes on Romans 4 at \hyperref[RIGHTEOUSNESS]{Righteousness}
			\end{compactitem}
		
	% -----
	\subsubsection{Thoughts from other people}
	% -----
		\label{FAITH-others}
		
		\begin{compactdesc}
			\item[27 December 2020, Roger Johnson] Having/exercising faith involves seeing the world in a certain way, seeing ourselves in a certain way. It may involve completely reworking how you understand the world because some people (I am one of these) are so accustomed to depending on themselves and doing things in a certain way, and exercising faith may require you not to depend on yourself, and may require you to change how you see things working in the world, or change how you assume they work. It usually means you will have to trust more.
		\end{compactdesc}